% 5) Indicadores quantitativos. Indicação separada das quantidades totais de:
% 1) livros publicados;
% 2) publicações em periódicos com seletiva política editorial;
% 3) capítulos de livros;
% 4) dissertações de Mestrado:
%   4.a) orientadas e já defendidas, e
%   4.b) em andamento;
% 5) teses de Doutorado:
%   5.a) orientadas e já defendidas, e
%   5.b) em andamento;
% 6) supervisões de Pós-Doutorado:
%   6.a) concluídas e,
%   6.b) em andamento;
% 7) quantidade de citações recebidas na literatura científica internacional, segundo o Web of Science, Scopus ou Google Scholar;
% 8) patentes solicitadas, concedidas e licenciadas;
% 9) produtos desenvolvidos e lançados no mercado;
% 10) processos otimizados implementados em empresas ou organizações sociais;
% 11) empresas criadas ou apoiadas;
% 12) consultorias técnicas e científicas relevantes.
% Importa o CSV indicadores.csv e gera a tabela
% Certifique-se de que o arquivo indicadores.csv está no mesmo diretório
% Dica: abra o indicadores.csv no Excel, edite e salve como "CSV UTF-8" (veja o README).

\pgfplotstabletypeset[
    col sep=semicolon, % Define que o separador é ponto e vírgula
    header=true, % A primeira linha do CSV é o cabeçalho (as colunas são lidas por posição, imune ao BOM do Excel)
    string type, % Define todas as colunas como texto
    every head row/.style={before row=\toprule,after row=\midrule}, % Estilo do cabeçalho
    every last row/.style={after row=\bottomrule}, % Estilo da última linha
    every even row/.style={before row={\rowcolor[gray]{0.85}}}, % Estilo das linhas pares
    every odd row/.style={before row={\rowcolor[gray]{0.95}}}, % Estilo das linhas ímpares
    display columns/0/.style={column name={\textbf{Indicador}}, column type={p{12cm}}}, % Define a largura e título da coluna Tipo
    display columns/1/.style={column name={\textbf{Quantidade}}, column type={c}}, % Define o título e alinhamento da coluna Número
    columns={[index]0,[index]1},
]{indicadores.csv}

\vspace{1em}

Para ver o \emph{Curriculum Lattes} do pesquisador, acesse \href{http://lattes.cnpq.br/\idLattes}{http://lattes.cnpq.br/\idLattes}.